% !PRINTED 396
\weberpage{396}
%<<CONT
$Q_n : c$ zu setzen und $t$ aus (5) oder auch, nachdem $u$ bestimmt
ist, direct aus (7) zu ermitteln.

Ist $\nu$ ungerade und $n$ ein ungerades Vielfaches von $\nu$, so ist
\weq{8}{t^2 - D u^2 = -4,}
ist aber $\nu$ gerade oder $n$ ein gerades Vielfaches von $\nu$, so ist
\weq{9}{t^2 - D u^2 = 4.}
Die Gleichung (9) hat die selbstverständliche Lösung $t = 2$,
$u = 0$; in allen anderen Lösungen von (8) oder (9) sind $t$ und
$u$ von Null verschieden und wir betrachten hier nur die positiven
Lösungen.

Wir wollen nun nachweisen, dass man aus einer beliebigen
Gleichung von der Form (3) durch die Formeln (4), (5) \gesperrt{alle
Lösungen} von (8) und (9) erhält.

Nehmen wir also an, es sei $t$, $u$ irgend eine positive Lösung
der Gleichung (8) oder (9), und $\omega = \{a, b, c\}$ eine zur Discriminante $D$ gehörige reducirte irrationale Zahl. Wir bestimmen
durch die Gleichungen (6) $P_n$, $Q_n$, $P_{n-1}$, $Q_{n-1}$ die offenbar ganze
Zahlen werden, da $t$ und $u b$ entweder beide gerade oder beide
ungerade sind. Aus der Gleichung (3), der $\omega$ genügt, ergiebt sich
dann (2), woraus wieder auf
\weq{10}{\omega = \dfrac{P_n \omega + P_{n-1}}{Q_n \omega + Q_{n-1}}}
zu schliessen ist; und aus (8) oder (9) ergiebt sich
\weq{11}{P_n Q_{n-1} - Q_n P_{n-1} = \mp 1,}
worin, wie auch in den folgenden Formeln, das obere oder das
untere Zeichen gilt, je nachdem $t$, $u$ die Gleichung (8) oder (9)
befriedigt. Nun folgt aber aus (6)
\[
Q_n - Q_{n-1} = \frac{(2 c + b) u - t}{2},
\]
und weil $\{a, b, c\}$ eine reducirte Zahl ist, so ist $2 c + b > \sqrt{D}$
[\S.\,125, (4)], also
\[
Q_n - Q_{n-1} > \frac{u \sqrt{D} - t}{2} = \frac{u^2 D - t^2}{2(u \sqrt{D} + t)} = \frac{\pm 2}{u \sqrt{D} + t};
\]
da aber $u$ und $t$ positive ganze Zahlen sind und $D > 1$, so ist
$u \sqrt{D} + t > 2$, also $Q_n - Q_{n-1} > 0$ für das obere und $> -1$
für das untere Zeichen.

Weil nun $Q_n - Q_{n-1}$ eine ganze Zahl sein muss, so ist
sie hiernach gleich oder grösser als Null, und wir haben
%CONT>>
